\documentclass[11pt,a4paper]{article}
\usepackage{fontspec,amsmath,amssymb,amsthm,microtype}
\setmainfont{lmroman10-regular.otf}[BoldFont=lmroman10-bold.otf,ItalicFont=lmroman10-italic.otf,BoldItalicFont=lmroman10-bolditalic.otf]
\usepackage{xcolor}
\usepackage[most]{tcolorbox}
\usepackage{mathrsfs}
\providecommand{\lt}{<}
\providecommand{\gt}{>}
\setmonofont{DejaVuSansMono.ttf}[BoldFont=DejaVuSansMono-Bold.ttf,ItalicFont=DejaVuSansMono-Oblique.ttf,BoldItalicFont=DejaVuSansMono-BoldOblique.ttf,Scale=0.86]
\newfontfamily\leanfallback{STIXTwoMath-Regular.otf}[Scale=0.86]
\definecolor{leanrule}{gray}{0.62}
\definecolor{leanback}{gray}{0.962}
\definecolor{leantitle}{gray}{0.22}
\definecolor{leanmuted}{gray}{0.38}
\newtcolorbox{leanbox}[1]{enhanced,breakable,colback=leanback,colframe=leanrule,
  boxrule=0.4pt,arc=0.6mm,left=2.4mm,right=2.4mm,top=1.8mm,bottom=1.8mm,
  before skip=0.8em,after skip=1em,fonttitle=\small\bfseries,coltitle=leantitle,
  colbacktitle=leanback,titlerule=0.3pt,title={#1}}
\newlength{\leancharwidth}
\newenvironment{leancode}{\par\smallskip\begingroup\ttfamily\footnotesize\raggedright
  \settowidth{\leancharwidth}{M}\setlength{\parskip}{0pt}\setlength{\parindent}{0pt}
  \frenchspacing}{\par\endgroup\smallskip}
\newcommand{\leanline}[2]{\hangindent=\dimexpr#1\leancharwidth+4\leancharwidth\relax
  \hangafter=1\noindent\hspace*{#1\leancharwidth}#2\par}
\newcommand{\leansource}[1]{\par{\footnotesize\color{leanmuted}#1}\par}
\newcommand{\leanlabel}[1]{\par\smallskip\noindent{\small\bfseries #1}\par}
\usepackage[margin=25mm]{geometry}
\usepackage[unicode,hidelinks]{hyperref}
\hypersetup{pdfauthor={LangevinLean Project}}
\setlength{\emergencystretch}{3em}
\setlength{\parskip}{0.4em}
\setlength{\parindent}{0pt}
\allowdisplaybreaks[2]
\raggedbottom
\title{\texorpdfstring{Independent copies and the empirical parameter-force residual}{Independent copies and the empirical parameter-force residual}}
\author{LangevinLean Project}
\date{}
\begin{document}\maketitle
Draft mathematical exposition; these notes may contain mistakes.


Independent copies of the mean-field state give an empirical drift residual with exact \(N^{-1/2}\) norm factor. Starting with a joint random state on \(Q\), use the finite product \(P=Q^{\otimes N}\) and its coordinate pullbacks. When the reference parameter is deterministic, all copies have the same parameter while their latent components are independent. The copy maps are bounded and linear, so they commute with the reference integral equation. All latent drift residuals cancel; only the centered empirical parameter force remains.


\section{\texorpdfstring{The measure-preserving coordinate copy map}{The measure-preserving coordinate copy map}}


For product coordinate \(\omega_i\) define \(C_iX=X\circ\omega_i\). Evaluation preserves the reference measure, so pullback is an \(L^2\) isometry and a bounded linear map. Its representative equals \(X(\omega_i)\) almost surely. Coordinate substitution gives \(C_i(CX)=C(C_iX)\) for bounded linear \(C\), and constants satisfy \(C_i(\operatorname{const}x)=\operatorname{const}x\).



The evaluation \(\omega\mapsto\omega_i\) has pushforward \(Q\); hence integration of \(\|X(\omega_i)\|^2\) gives the original second moment. Pullback is consequently an isometry on \(L^2\). These pointwise and integral identities justify copying the projected reference equations.

\begin{leanbox}{Lean: \texttt{iidCopyCLM}, \texttt{iidCopyCLM\_\allowbreak{}coe}, \texttt{iidCopyCLM\_\allowbreak{}compLp} and 1 more}
\leansource{\texttt{iidCopyCLM} \quad ProximalLangevin/\allowbreak{}MeanFieldCopies.lean, lines 26–29}
\begin{leancode}
\leanline{0}{def\ iidCopyCLM\ (i\ :\ Fin\ N)\ :}
\leanline{4}{Lp\ E\ 2\ Q\ →L[ℝ]\ Lp\ E\ 2\ (Measure.pi\ (fun\ \_\ :\ Fin\ N\ =>\ Q))\ :=}
\leanline{2}{(Lp.compMeasurePreservingₗᵢ\ ℝ\ (fun\ ω\ :\ Fin\ N\ →\ Ω\ =>\ ω\ i)}
\leanline{4}{(measurePreserving\_eval\ (fun\ \_\ :\ Fin\ N\ =>\ Q)\ i)).toContinuousLinearMap}
\end{leancode}
\leansource{\texttt{iidCopyCLM\_\allowbreak{}coe} \quad ProximalLangevin/\allowbreak{}MeanFieldCopies.lean, lines 31–33}
\begin{leancode}
\leanline{0}{theorem\ iidCopyCLM\_coe\ (i\ :\ Fin\ N)\ (X\ :\ Lp\ E\ 2\ Q)\ :}
\leanline{4}{iidCopyCLM\ Q\ i\ X\ =ᵐ[Measure.pi\ (fun\ \_\ :\ Fin\ N\ =>\ Q)]\ fun\ ω\ =>\ X\ (ω\ i)}
\end{leancode}
\leansource{\texttt{iidCopyCLM\_\allowbreak{}compLp} \quad ProximalLangevin/\allowbreak{}MeanFieldCopies.lean, lines 35–42}
\begin{leancode}
\leanline{0}{theorem\ iidCopyCLM\_compLp\ (i\ :\ Fin\ N)\ (C\ :\ E\ →L[ℝ]\ F)\ (X\ :\ Lp\ E\ 2\ Q)\ :}
\leanline{4}{iidCopyCLM\ Q\ i\ (C.compLp\ X)\ =\ C.compLp\ (iidCopyCLM\ Q\ i\ X)}
\end{leancode}
\leansource{\texttt{iidCopyCLM\_\allowbreak{}const} \quad ProximalLangevin/\allowbreak{}MeanFieldCopies.lean, lines 44–52}
\begin{leancode}
\leanline{0}{theorem\ iidCopyCLM\_const\ (i\ :\ Fin\ N)\ (x\ :\ E)\ :}
\leanline{4}{iidCopyCLM\ Q\ i\ (Lp.const\ 2\ Q\ x)=Lp.const\ 2\ (Measure.pi\ (fun\ \_\ :\ Fin\ N\ =>\ Q))\ x}
\end{leancode}
\leanlabel{How the Lean proof follows the argument}
Formal correspondence: iidCopyCLM : A genuine iid copy of a random variable on the finite product space.



\texttt{iid\allowbreak{}Copy\allowbreak{}CLM} is a definition with defining expression \texttt{(Lp.\allowbreak{}comp\allowbreak{}Measure\allowbreak{}Preserving\ensuremath{{}_{li}} \ensuremath{\mathbb R} (fun \ensuremath{\omega} : Fin N \ensuremath{\to} \ensuremath{\Omega} => \ensuremath{\omega} i)}. The right-hand side specifies the object used in the mathematical statement.



\texttt{iid\allowbreak{}Copy\allowbreak{}CLM\_\allowbreak{}coe}: the proof is the application \texttt{Lp.\allowbreak{}coe\allowbreak{}Fn\_\allowbreak{}comp\allowbreak{}Measure\allowbreak{}Preserving X (measure\allowbreak{}Preserving\_\allowbreak{}eval (fun \_\allowbreak{} : Fin N => Q) i)}. It specializes that result with the hypotheses and objects in the exact statement.



\texttt{iid\allowbreak{}Copy\allowbreak{}CLM\_\allowbreak{}comp\allowbreak{}Lp}: the intermediate fact is \texttt{hp}, in source order; the source steps apply or rewrite \texttt{Lp.\allowbreak{}ext}.



\texttt{iid\allowbreak{}Copy\allowbreak{}CLM\_\allowbreak{}const}: the intermediate fact is \texttt{hp}, in source order; the source steps apply or rewrite \texttt{Lp.\allowbreak{}ext}.
\end{leanbox}

\section{\texorpdfstring{The centered empirical fluctuation norm}{The centered empirical fluctuation norm}}


For \(N>0\) define \(\mathcal C X=X-\operatorname{const}(\mathbb EX)\) and \(\Delta_N(X)=N^{-1}\sum_iC_iX-\operatorname{const}(\mathbb EX)\). Finite-sum representative identities identify the latter with the centered sample mean. Independence on the product space gives

\[\|\Delta_N(X)\|_{L^2(P)}^2=\frac1N\|\mathcal C X\|_{L^2(Q)}^2,\qquad\|\Delta_N(X)\|_2=\frac{\|\mathcal C X\|_2}{\sqrt N}.\]

The pullback \(X\mapsto X(\omega_i)\) preserves \(Q\) and is a linear isometry. Thus the centered copy sum is an \(L^2\) vector represented by that sample mean. The exact variance statement is

\[\|\text{empirical fluctuation}\|_2^2=\frac{\|X-\mathbb EX\|_2^2}{N},\]

Taking nonnegative square roots gives factor \(N^{-1/2}\). The centered force norm is divided by \(\sqrt N\). This uses independence of reference copies; the later interacting particle states are not assumed independent.

\begin{leanbox}{Lean: \texttt{lpCenterCLM}, \texttt{iidEmpiricalFluctuation}, \texttt{iidEmpiricalFluctuation\_\allowbreak{}coe} and 1 more}
\leansource{\texttt{lpCenterCLM} \quad ProximalLangevin/\allowbreak{}MeanFieldCopies.lean, lines 62–63}
\begin{leancode}
\leanline{0}{def\ lpCenterCLM\ :\ Lp\ E\ 2\ Q\ →L[ℝ]\ Lp\ E\ 2\ Q\ :=}
\leanline{2}{ContinuousLinearMap.id\ ℝ\ \_-(Lp.constL\ 2\ Q\ ℝ).comp\ (lpExpectationCLM\ Q)}
\end{leancode}
\leansource{\texttt{iidEmpiricalFluctuation} \quad ProximalLangevin/\allowbreak{}MeanFieldCopies.lean, lines 66–69}
\begin{leancode}
\leanline{0}{def\ iidEmpiricalFluctuation\ (X\ :\ Lp\ E\ 2\ Q)\ :}
\leanline{4}{Lp\ E\ 2\ (Measure.pi\ (fun\ \_\ :\ Fin\ N\ =>\ Q))\ :=}
\leanline{2}{(N\ :\ ℝ)⁻¹\ •\ ∑\ i,iidCopyCLM\ Q\ i\ X-}
\leanline{4}{Lp.const\ 2\ (Measure.pi\ (fun\ \_\ :\ Fin\ N\ =>\ Q))\ (lpExpectationCLM\ Q\ X)}
\end{leancode}
\leansource{\texttt{iidEmpiricalFluctuation\_\allowbreak{}coe} \quad ProximalLangevin/\allowbreak{}MeanFieldCopies.lean, lines 72–84}
\begin{leancode}
\leanline{0}{theorem\ iidEmpiricalFluctuation\_coe\ (X\ :\ Lp\ E\ 2\ Q)\ :}
\leanline{4}{iidEmpiricalFluctuation\ Q\ (N\ :=\ N)\ X\ =ᵐ[Measure.pi\ (fun\ \_\ :\ Fin\ N\ =>\ Q)]}
\leanline{6}{fun\ ω\ =>\ empiricalAverage\ X\ ω-(∫\ a,X\ a\ ∂Q)}
\end{leancode}
\leansource{\texttt{iidEmpiricalFluctuation\_\allowbreak{}norm} \quad ProximalLangevin/\allowbreak{}MeanFieldCopies.lean, lines 86–106}
\begin{leancode}
\leanline{0}{theorem\ iidEmpiricalFluctuation\_norm\ (hN\ :\ 0<N)\ (X\ :\ Lp\ E\ 2\ Q)\ :}
\leanline{4}{‖iidEmpiricalFluctuation\ Q\ (N\ :=\ N)\ X‖=‖lpCenterCLM\ Q\ X‖/Real.sqrt\ (N\ :\ ℝ)}
\end{leancode}
\leanlabel{How the Lean proof follows the argument}
Formal correspondence: lpCenterCLM : Center an actual square-integrable random variable by its expectation. iidEmpiricalFluctuation : The actual finite-product centered empirical mean as an \(L^2\) random variable.



\texttt{lp\allowbreak{}Center\allowbreak{}CLM} is a definition with defining expression \texttt{Continuous\allowbreak{}Linear\allowbreak{}Map.\allowbreak{}id \ensuremath{\mathbb R} \_\allowbreak{}-\allowbreak{}(Lp.\allowbreak{}const\allowbreak{}L 2 Q \ensuremath{\mathbb R}).\allowbreak{}comp (lp\allowbreak{}Expectation\allowbreak{}CLM Q)}. The right-hand side specifies the object used in the mathematical statement.



\texttt{iid\allowbreak{}Empirical\allowbreak{}Fluctuation} is a definition with defining expression \texttt{(N : \ensuremath{\mathbb R})\ensuremath{{}^{-1}} \ensuremath{\bullet} ∑ i,iid\allowbreak{}Copy\allowbreak{}CLM Q i X-\allowbreak{}}. The right-hand side specifies the object used in the mathematical statement.



\texttt{iid\allowbreak{}Empirical\allowbreak{}Fluctuation\_\allowbreak{}coe}: the intermediate fact is \texttt{hi}, in source order; the source steps apply or rewrite \texttt{Pi.\allowbreak{}sub\_\allowbreak{}apply}, \texttt{Pi.\allowbreak{}smul\_\allowbreak{}apply}, \texttt{lp\allowbreak{}Expectation\allowbreak{}CLM\_\allowbreak{}apply}, \texttt{Function.\allowbreak{}const\_\allowbreak{}apply}.



\texttt{iid\allowbreak{}Empirical\allowbreak{}Fluctuation\_\allowbreak{}norm}: the intermediate facts include \texttt{hn}, \texttt{hv}, \texttt{hf}, \texttt{heq}, \texttt{hs}, in source order; the source steps apply or rewrite \texttt{sub\_\allowbreak{}apply}, \texttt{Continuous\allowbreak{}Linear\allowbreak{}Map.\allowbreak{}id\_\allowbreak{}apply}, \texttt{norm\_\allowbreak{}nonneg}, \texttt{Real.\allowbreak{}sqrt\_\allowbreak{}div}, \texttt{sq\_\allowbreak{}nonneg}.
\end{leanbox}

\section{\texorpdfstring{The reference particle state with deterministic parameter}{The reference particle state with deterministic parameter}}


Define the reference particle state by

\[\mathcal R_N(X)=\left(\operatorname{const}(\mathbb E P_\theta X),\ (P_xC_iX)_i\right).\]

Bounded linear maps commute with expectation, and a constant has its original value as expectation under a probability measure. Thus if \(P_\theta X=\operatorname{const}\theta\), the parameter component of \(\mathcal R_N(X)\) is \(\operatorname{const}\theta\).



The latent coordinates are independent coordinate copies, while the deterministic parameter is unchanged. The construction is a bounded linear map, which permits passage through the reference time integral.

\begin{leanbox}{Lean: \texttt{lpExpectationCLM\_\allowbreak{}compLp}, \texttt{lpExpectationCLM\_\allowbreak{}const}, \texttt{iidMeanFieldCopiesCLM} and 3 more}
\leansource{\texttt{lpExpectationCLM\_\allowbreak{}compLp} \quad ProximalLangevin/\allowbreak{}MeanFieldCopies.lean, lines 115–120}
\begin{leancode}
\leanline{0}{theorem\ lpExpectationCLM\_compLp\ \{E\ F\ :\ Type*\}\ [NormedAddCommGroup\ E]\ [NormedSpace\ ℝ\ E]}
\leanline{4}{[CompleteSpace\ E]\ [NormedAddCommGroup\ F]\ [NormedSpace\ ℝ\ F]\ [CompleteSpace\ F]}
\leanline{4}{(C\ :\ E\ →L[ℝ]\ F)\ (X\ :\ Lp\ E\ 2\ Q)\ :}
\leanline{4}{lpExpectationCLM\ Q\ (C.compLp\ X)=C\ (lpExpectationCLM\ Q\ X)}
\end{leancode}
\leansource{\texttt{lpExpectationCLM\_\allowbreak{}const} \quad ProximalLangevin/\allowbreak{}MeanFieldCopies.lean, lines 122–125}
\begin{leancode}
\leanline{0}{theorem\ lpExpectationCLM\_const\ \{E\ :\ Type*\}\ [NormedAddCommGroup\ E]\ [NormedSpace\ ℝ\ E]}
\leanline{4}{[CompleteSpace\ E]\ (x\ :\ E)\ :\ lpExpectationCLM\ Q\ (Lp.const\ 2\ Q\ x)=x}
\end{leancode}
\leansource{\texttt{iidMeanFieldCopiesCLM} \quad ProximalLangevin/\allowbreak{}MeanFieldCopies.lean, lines 128–133}
\begin{leancode}
\leanline{0}{def\ iidMeanFieldCopiesCLM\ :\ Lp\ (IPLA.Joint\ dθ\ dx)\ 2\ Q\ →L[ℝ]}
\leanline{4}{MeanFieldParticleState\ (Measure.pi\ (fun\ \_\ :\ Fin\ N\ =>\ Q))\ N\ dθ\ dx\ :=}
\leanline{2}{((Lp.constL\ 2\ (Measure.pi\ (fun\ \_\ :\ Fin\ N\ =>\ Q))\ ℝ).comp}
\leanline{4}{(jointParameterCLM.comp\ (lpExpectationCLM\ Q))).prod}
\leanline{4}{(ContinuousLinearMap.pi\ (fun\ i\ :\ Fin\ N\ =>}
\leanline{6}{((@jointLatentCLM\ dθ\ dx).compLpL\ 2\ (Measure.pi\ (fun\ \_\ :\ Fin\ N\ =>\ Q))).comp\ (iidCopyCLM\ Q\ i)))}
\end{leancode}
\leansource{\texttt{iidMeanFieldCopiesCLM\_\allowbreak{}parameter} \quad ProximalLangevin/\allowbreak{}MeanFieldCopies.lean, lines 135–137}
\begin{leancode}
\leanline{0}{@[simp]\ theorem\ iidMeanFieldCopiesCLM\_parameter\ (X\ :\ Lp\ (IPLA.Joint\ dθ\ dx)\ 2\ Q)\ :}
\leanline{4}{(iidMeanFieldCopiesCLM\ Q\ (N\ :=\ N)\ X).1=}
\leanline{6}{Lp.const\ 2\ (Measure.pi\ (fun\ \_\ :\ Fin\ N\ =>\ Q))\ (jointParameterCLM\ (lpExpectationCLM\ Q\ X))}
\end{leancode}
\leansource{\texttt{iidMeanFieldCopiesCLM\_\allowbreak{}latent} \quad ProximalLangevin/\allowbreak{}MeanFieldCopies.lean, lines 139–140}
\begin{leancode}
\leanline{0}{@[simp]\ theorem\ iidMeanFieldCopiesCLM\_latent\ (X\ :\ Lp\ (IPLA.Joint\ dθ\ dx)\ 2\ Q)\ (i\ :\ Fin\ N)\ :}
\leanline{4}{(iidMeanFieldCopiesCLM\ Q\ (N\ :=\ N)\ X).2\ i=jointLatentCLM.compLp\ (iidCopyCLM\ Q\ i\ X)}
\end{leancode}
\leansource{\texttt{iidMeanFieldCopiesCLM\_\allowbreak{}parameter\_\allowbreak{}const} \quad ProximalLangevin/\allowbreak{}MeanFieldCopies.lean, lines 142–148}
\begin{leancode}
\leanline{0}{theorem\ iidMeanFieldCopiesCLM\_parameter\_const\ (X\ :\ Lp\ (IPLA.Joint\ dθ\ dx)\ 2\ Q)}
\leanline{4}{(θ\ :\ EuclideanSpace\ ℝ\ (Fin\ dθ))\ (hθ\ :\ jointParameterCLM.compLp\ X=Lp.const\ 2\ Q\ θ)\ :}
\leanline{4}{(iidMeanFieldCopiesCLM\ Q\ (N\ :=\ N)\ X).1=Lp.const\ 2\ (Measure.pi\ (fun\ \_\ :\ Fin\ N\ =>\ Q))\ θ}
\end{leancode}
\leanlabel{How the Lean proof follows the argument}
Formal correspondence: lpExpectationCLM\_compLp : Actual expectation commutes with any bounded linear map. iidMeanFieldCopiesCLM : Mean-field parameter expectation and genuine iid latent copies.



\texttt{lp\allowbreak{}Expectation\allowbreak{}CLM\_\allowbreak{}comp\allowbreak{}Lp}: the source steps apply or rewrite \texttt{lp\allowbreak{}Expectation\allowbreak{}CLM\_\allowbreak{}apply}, \texttt{integral\_\allowbreak{}congr\_\allowbreak{}ae}, \texttt{C.\allowbreak{}coe\allowbreak{}Fn\_\allowbreak{}comp\allowbreak{}Lp}, \texttt{C.\allowbreak{}integral\_\allowbreak{}comp\_\allowbreak{}comm}, \texttt{Lp.\allowbreak{}mem\allowbreak{}Lp}.



\texttt{lp\allowbreak{}Expectation\allowbreak{}CLM\_\allowbreak{}const}: the source steps apply or rewrite \texttt{lp\allowbreak{}Expectation\allowbreak{}CLM\_\allowbreak{}apply}, \texttt{integral\_\allowbreak{}congr\_\allowbreak{}ae}, \texttt{integral\_\allowbreak{}const}, \texttt{prob\allowbreak{}Real\_\allowbreak{}univ}, \texttt{one\_\allowbreak{}smul}.



\texttt{iid\allowbreak{}Mean\allowbreak{}Field\allowbreak{}Copies\allowbreak{}CLM} is a definition with defining expression \texttt{((Lp.\allowbreak{}const\allowbreak{}L 2 (Measure.\allowbreak{}pi (fun \_\allowbreak{} : Fin N => Q)) \ensuremath{\mathbb R}).\allowbreak{}comp}. The right-hand side specifies the object used in the mathematical statement.



\texttt{iid\allowbreak{}Mean\allowbreak{}Field\allowbreak{}Copies\allowbreak{}CLM\_\allowbreak{}parameter} is proved by \texttt{rfl}: evaluating the definitions makes the two sides the same expression.



\texttt{iid\allowbreak{}Mean\allowbreak{}Field\allowbreak{}Copies\allowbreak{}CLM\_\allowbreak{}latent} is proved by \texttt{rfl}: evaluating the definitions makes the two sides the same expression.



\texttt{iid\allowbreak{}Mean\allowbreak{}Field\allowbreak{}Copies\allowbreak{}CLM\_\allowbreak{}parameter\_\allowbreak{}const}: the intermediate fact is \texttt{h}, in source order; the source steps apply or rewrite \texttt{iid\allowbreak{}Mean\allowbreak{}Field\allowbreak{}Copies\allowbreak{}CLM\_\allowbreak{}parameter}, \texttt{lp\allowbreak{}Expectation\allowbreak{}CLM\_\allowbreak{}comp\allowbreak{}Lp}, \texttt{lp\allowbreak{}Expectation\allowbreak{}CLM\_\allowbreak{}const}.
\end{leanbox}


Earlier results used here: \cite[Section~3.1.1]{LangevinLean2026IPLA}; \cite[Section~3.11.1]{LangevinLean2026IPLA}.

\section{\texorpdfstring{Reconstruction and drift composition for copied pairs}{Reconstruction and drift composition for copied pairs}}


Recombining the parameter and latent projections recovers the original \(L^2\) vector. For deterministic parameter this gives \(J_i\mathcal R_N(X)=C_iX\). Composition of the copied state with a Lipschitz pointwise drift is the same as copying the composition, since both representatives equal \(b(X(\omega_i))\) almost surely.



For the gradient, this common representative is \(\nabla V(X(\omega_i))\). Hence each reconstructed reference pair has precisely the copied joint drift. This equality causes the latent components to cancel in the particle-minus-reference drift.

\begin{leanbox}{Lean: \texttt{lpJointPair\_\allowbreak{}projections}, \texttt{iidMeanFieldCopies\_\allowbreak{}joint}, \texttt{iidCopyCLM\_\allowbreak{}lpDrift}}
\leansource{\texttt{lpJointPair\_\allowbreak{}projections} \quad ProximalLangevin/\allowbreak{}MeanFieldCopies.lean, lines 151–159}
\begin{leancode}
\leanline{0}{theorem\ lpJointPair\_projections\ \{A\ :\ Type*\}\ [MeasurableSpace\ A]\ (P\ :\ Measure\ A)}
\leanline{4}{[IsProbabilityMeasure\ P]\ (X\ :\ Lp\ (IPLA.Joint\ dθ\ dx)\ 2\ P)\ :}
\leanline{4}{lpJointPair\ P\ (jointParameterCLM.compLp\ X)\ (jointLatentCLM.compLp\ X)=X}
\end{leancode}
\leansource{\texttt{iidMeanFieldCopies\_\allowbreak{}joint} \quad ProximalLangevin/\allowbreak{}MeanFieldCopies.lean, lines 163–173}
\begin{leancode}
\leanline{0}{theorem\ iidMeanFieldCopies\_joint\ (X\ :\ Lp\ (IPLA.Joint\ dθ\ dx)\ 2\ Q)}
\leanline{4}{(θ\ :\ EuclideanSpace\ ℝ\ (Fin\ dθ))\ (hθ\ :\ jointParameterCLM.compLp\ X=Lp.const\ 2\ Q\ θ)}
\leanline{4}{(i\ :\ Fin\ N)\ :}
\leanline{4}{lpJointPair\ (Measure.pi\ (fun\ \_\ :\ Fin\ N\ =>\ Q))}
\leanline{6}{(iidMeanFieldCopiesCLM\ Q\ (N\ :=\ N)\ X).1\ ((iidMeanFieldCopiesCLM\ Q\ X).2\ i)=iidCopyCLM\ Q\ i\ X}
\end{leancode}
\leansource{\texttt{iidCopyCLM\_\allowbreak{}lpDrift} \quad ProximalLangevin/\allowbreak{}MeanFieldCopies.lean, lines 176–185}
\begin{leancode}
\leanline{0}{theorem\ iidCopyCLM\_lpDrift\ \{E\ :\ Type*\}\ [NormedAddCommGroup\ E]\ [NormedSpace\ ℝ\ E]\ (g\ :\ E\ →\ E)\ \{L\ :\ ℝ≥0\}}
\leanline{4}{(hg\ :\ LipschitzWith\ L\ g)\ (X\ :\ Lp\ E\ 2\ Q)\ (i\ :\ Fin\ N)\ :}
\leanline{4}{ULA.lpDrift\ (Measure.pi\ (fun\ \_\ :\ Fin\ N\ =>\ Q))\ g\ hg\ (iidCopyCLM\ Q\ i\ X)=}
\leanline{6}{iidCopyCLM\ Q\ i\ (ULA.lpDrift\ Q\ g\ hg\ X)}
\end{leancode}
\leanlabel{How the Lean proof follows the argument}
Formal correspondence: lpJointPair\_projections : Reconstructing one pair from a joint random variable loses no coordinates. iidMeanFieldCopies\_joint : Each reference pair is the selected genuine joint copy when its parameter is deterministic. This is an equality of actual \(L^2\) random variables. iidCopyCLM\_lpDrift : Actual pointwise Lipschitz drift commutes with an iid pullback.



\texttt{lp\allowbreak{}Joint\allowbreak{}Pair\_\allowbreak{}projections}: the source steps apply or rewrite \texttt{Lp.\allowbreak{}ext}.



\texttt{iid\allowbreak{}Mean\allowbreak{}Field\allowbreak{}Copies\_\allowbreak{}joint}: the intermediate fact is \texttt{hp}, in source order; the source steps apply or rewrite \texttt{iid\allowbreak{}Mean\allowbreak{}Field\allowbreak{}Copies\allowbreak{}CLM\_\allowbreak{}parameter\_\allowbreak{}const}, \texttt{iid\allowbreak{}Mean\allowbreak{}Field\allowbreak{}Copies\allowbreak{}CLM\_\allowbreak{}latent}, \texttt{iid\allowbreak{}Copy\allowbreak{}CLM\_\allowbreak{}comp\allowbreak{}Lp}, \texttt{iid\allowbreak{}Copy\allowbreak{}CLM\_\allowbreak{}const}, \texttt{lp\allowbreak{}Joint\allowbreak{}Pair\_\allowbreak{}projections}.



\texttt{iid\allowbreak{}Copy\allowbreak{}CLM\_\allowbreak{}lp\allowbreak{}Drift}: the intermediate fact is \texttt{hsrc}, in source order; the source steps apply or rewrite \texttt{Lp.\allowbreak{}ext}.
\end{leanbox}


Earlier results used here: \cite[Section~3.13.2]{LangevinLean2026ULASC}; \cite[Section~3.13.3]{LangevinLean2026ULASC}; \cite[Section~3.1.1]{LangevinLean2026IPLA}.

\section{\texorpdfstring{The empirical parameter-force drift residual}{The empirical parameter-force drift residual}}


Apply the particle drift to \(\mathcal R_N(X)\) and compare it with \(\mathcal R_N(\mathcal A_V(X))\). Reconstruction and gradient commutation identify all latent components, leaving

\[\mathcal A_N(\mathcal R_NX)-\mathcal R_N(\mathcal A_VX)=\left(\Delta_N(P_\theta\nabla V(X)),0\right).\]

The maximum product norm and sample-mean variance identity give residual norm \(\|\mathcal C(P_\theta\nabla V(X))\|_2/\sqrt N\). In detail, latent coordinate \(i\) combines with the deterministic parameter to recover joint copy \(i\). The parameter discrepancy is

\[\frac1N\sum_i\nabla_\theta V(\theta,X_i)-\mathbb E\nabla_\theta V(\theta,X).\]

It is the empirical parameter force minus its expectation. Its norm is the centered single-copy force norm divided by \(\sqrt N\), equivalently \(\|\mathcal C(P_\theta\nabla V(X))\|_2/\sqrt N\). This exact residual supplies the sampling forcing used in the comparison theorem.

\begin{leanbox}{Lean: \texttt{iidMeanFieldCopies\_\allowbreak{}drift\_\allowbreak{}residual}, \texttt{iidMeanFieldCopies\_\allowbreak{}drift\_\allowbreak{}residual\_\allowbreak{}norm}}
\leansource{\texttt{iidMeanFieldCopies\_\allowbreak{}drift\_\allowbreak{}residual} \quad ProximalLangevin/\allowbreak{}MeanFieldCopies.lean, lines 189–215}
\begin{leancode}
\leanline{0}{theorem\ iidMeanFieldCopies\_drift\_residual\ (V\ :\ IPLA.Joint\ dθ\ dx\ →\ ℝ)\ \{L\ :\ ℝ≥0\}}
\leanline{4}{(hLip\ :\ LipschitzWith\ L\ (gradient\ V))\ (X\ :\ Lp\ (IPLA.Joint\ dθ\ dx)\ 2\ Q)}
\leanline{4}{(θ\ :\ EuclideanSpace\ ℝ\ (Fin\ dθ))\ (hθ\ :\ jointParameterCLM.compLp\ X=Lp.const\ 2\ Q\ θ)\ :}
\leanline{4}{lpParticleDrift\ (Measure.pi\ (fun\ \_\ :\ Fin\ N\ =>\ Q))\ V\ hLip\ (iidMeanFieldCopiesCLM\ Q\ (N\ :=\ N)\ X)-}
\leanline{6}{iidMeanFieldCopiesCLM\ Q\ (lpMeanFieldDrift\ Q\ V\ hLip\ X)=}
\leanline{6}{(iidEmpiricalFluctuation\ Q\ (N\ :=\ N)}
\leanline{8}{(jointParameterCLM.compLp\ (ULA.lpDrift\ Q\ (gradient\ V)\ hLip\ X)),fun\ \_\ =>\ 0)}
\end{leancode}
\leansource{\texttt{iidMeanFieldCopies\_\allowbreak{}drift\_\allowbreak{}residual\_\allowbreak{}norm} \quad ProximalLangevin/\allowbreak{}MeanFieldCopies.lean, lines 218–231}
\begin{leancode}
\leanline{0}{theorem\ iidMeanFieldCopies\_drift\_residual\_norm\ (hN\ :\ 0<N)}
\leanline{4}{(V\ :\ IPLA.Joint\ dθ\ dx\ →\ ℝ)\ \{L\ :\ ℝ≥0\}\ (hLip\ :\ LipschitzWith\ L\ (gradient\ V))}
\leanline{4}{(X\ :\ Lp\ (IPLA.Joint\ dθ\ dx)\ 2\ Q)\ (θ\ :\ EuclideanSpace\ ℝ\ (Fin\ dθ))}
\leanline{4}{(hθ\ :\ jointParameterCLM.compLp\ X=Lp.const\ 2\ Q\ θ)\ :}
\leanline{4}{‖lpParticleDrift\ (Measure.pi\ (fun\ \_\ :\ Fin\ N\ =>\ Q))\ V\ hLip\ (iidMeanFieldCopiesCLM\ Q\ (N\ :=\ N)\ X)-}
\leanline{6}{iidMeanFieldCopiesCLM\ Q\ (lpMeanFieldDrift\ Q\ V\ hLip\ X)‖=}
\leanline{6}{‖lpCenterCLM\ Q\ (jointParameterCLM.compLp\ (ULA.lpDrift\ Q\ (gradient\ V)\ hLip\ X))‖/}
\leanline{8}{Real.sqrt\ (N\ :\ ℝ)}
\end{leancode}
\leanlabel{How the Lean proof follows the argument}
Formal correspondence: iidMeanFieldCopies\_drift\_residual : The finite-particle drift defect at the constructed iid reference is exactly its centered empirical parameter-gradient error. The latent defect is zero. iidMeanFieldCopies\_drift\_residual\_norm : The actual drift defect has its exact independent-sampling norm.



\texttt{iid\allowbreak{}Mean\allowbreak{}Field\allowbreak{}Copies\_\allowbreak{}drift\_\allowbreak{}residual}: the intermediate facts include \texttt{G}, \texttt{hg}, in source order; the source steps apply or rewrite \texttt{iid\allowbreak{}Mean\allowbreak{}Field\allowbreak{}Copies\_\allowbreak{}joint}, \texttt{iid\allowbreak{}Copy\allowbreak{}CLM\_\allowbreak{}lp\allowbreak{}Drift}, \texttt{iid\allowbreak{}Mean\allowbreak{}Field\allowbreak{}Copies\allowbreak{}CLM\_\allowbreak{}latent}, \texttt{iid\allowbreak{}Copy\allowbreak{}CLM\_\allowbreak{}comp\allowbreak{}Lp}, \texttt{sub\_\allowbreak{}self}.



\texttt{iid\allowbreak{}Mean\allowbreak{}Field\allowbreak{}Copies\_\allowbreak{}drift\_\allowbreak{}residual\_\allowbreak{}norm}: the source steps apply or rewrite \texttt{iid\allowbreak{}Mean\allowbreak{}Field\allowbreak{}Copies\_\allowbreak{}drift\_\allowbreak{}residual}, \texttt{Prod.\allowbreak{}norm\_\allowbreak{}mk}, \texttt{max\_\allowbreak{}eq\_\allowbreak{}left}, \texttt{norm\_\allowbreak{}nonneg}, \texttt{iid\allowbreak{}Empirical\allowbreak{}Fluctuation\_\allowbreak{}norm}.
\end{leanbox}


Earlier results used here: \cite[Section~3.13.2]{LangevinLean2026ULASC}; \cite[Section~3.1.1]{LangevinLean2026IPLA}.

\begin{thebibliography}{ULA-LSI}
\bibitem[ULA-SC]{LangevinLean2026ULASC}
LangevinLean Project. \emph{Unadjusted Langevin under Strong Convexity}. LL-01, 2026. Draft report, content revision d50c0f64815b. \url{https://odakyildiz.github.io/LangevinLean/docs/pdf/UnadjustedLangevin/FullProof.pdf}.
\bibitem[IPLA-SC]{LangevinLean2026IPLA}
LangevinLean Project. \emph{Interacting Particle Langevin under Strong Convexity}. LL-03, 2026. Draft report, content revision bfbb88fd1d79. \url{https://odakyildiz.github.io/LangevinLean/docs/pdf/IPLA/FullProof.pdf}.
\end{thebibliography}

\end{document}
